\documentclass[10pt,a4paper]{amsart}
\usepackage[margin=19mm]{geometry}
\usepackage{amsmath,amssymb,mathtools}
\usepackage[scr=rsfs,cal=euler]{mathalfa}
\usepackage{xfrac}
\usepackage[unicode,hidelinks,bookmarks=false]{hyperref}
\newcommand{\ie}{i.e.\ }
\newcommand{\cf}{cf.\ }
\newcommand{\resp}{respectively\ }
\newcommand{\Subsection}{\subsection}
\providecommand{\nobreakdash}{\nobreak\hbox{-}}
\setlength{\emergencystretch}{2em}
\multlinegap0pt
\usepackage{bm}
\usepackage{bbm}
\usepackage{dsfont}
\usepackage{xspace}
\usepackage{cancel}
\usepackage{mathtools}
\usepackage{amsthm}
\makeatletter
\@ifundefined{theorem}{\newtheorem{theorem}{Theorem}}{}
\@ifundefined{lemma}{\newtheorem{lemma}[theorem]{Lemma}}{}
\makeatother
\makeatletter
\expandafter\ifx\csname claimproof\endcsname\relax
\newenvironment{claimproof}{\begin{proof}}{\renewcommand{\qedsymbol}{$\dashv$}\end{proof}}
\fi
\makeatother
\title[Polymorphic Ordinal Notations]{Polymorphic Ordinal Notations}
\author{Henry Towsner}
\date{Source: arXiv:2504.02131v2 (CC BY 4.0)}
\begin{document}
\maketitle

\noindent\emph{Source and licence.} Polymorphic Ordinal Notations. Version arXiv:2504.02131v2 (CC BY 4.0), distributed under the Creative Commons Attribution 4.0 International licence, as recorded in the arXiv licence metadata for this exact version and shown on the abstract page https://arxiv.org/abs/2504.02131v2. Full rights notice: licenses/arxiv-2504.02131-CC-BY-4.0.txt.\par\smallskip

\noindent\textit{Editorial scope.} Counted unit: the Lemma whose source label is \texttt{None} in the article (source0.tex lines 409--411, source label \texttt{None}), delivered together with the article's own complete original proof of it (source0.tex lines 412--426, one \texttt{proof} environment of 15 lines). No mathematics is written, repaired or re-derived: statement and proof are the article's own text line for line. Required context retained uncounted: none. Every definition, display and label of those lines is retained unchanged. Omitted from this package and not redistributed: 1--408, 427--1023. Nothing inside a retained excerpt is omitted or altered: every retained body is delivered exactly as the article's lines are written, so no omission receipt of any kind accompanies this package. Citations: the retained excerpts carry no citation command at all, so no bibliography entry of the article is transcribed and no reference is redistributed. Reference adaptation: none was needed and none was made; every reference inside the delivered unit resolves inside it, so no unresolved reference or citation is produced. Printed slips kept literal: no printed word, symbol, hypothesis or number of the counted statement or of its proof was altered. Layout and class: the article's own class is \texttt{amsart} with packages including fontenc, inputenc, textcomp, lmodern, microtype, xcolor, amssymb, amsthm, amscd, mathabx, hyperref, mathscinet, biblatex, bussproofs; the wrapper is the lane's pinned \texttt{amsart} editorial wrapper at 10pt on A4, which restates the theorem-like environments the retained text needs and adds the article's own macro definitions that the retained text needs (none), with extra wrapper packages (bm, bbm, dsfont, xspace, cancel, mathtools). The delivered numbering is the unit's own. Assets: no retained span contains a figure environment, a graphics-inclusion command, a PDF-page inclusion, a table or a TikZ picture; no image file is copied into this package, the package declares no asset dependency and no image file is redistributed with it. Licence: version arXiv:2504.02131v2 of this article is distributed under the Creative Commons Attribution 4.0 International licence, as recorded in the arXiv licence metadata for this exact version and shown on the abstract page https://arxiv.org/abs/2504.02131v2. A full rights notice is delivered at licenses/arxiv-2504.02131-CC-BY-4.0.txt. Characters: every retained excerpt is pure ASCII, so the delivered engine, XeLaTeX with its default Latin Modern text font, reports no missing character.
\begin{lemma}
  If $\alpha\in\mathrm{Acc}_n$ and $n>-\infty$ then $(\vartheta\alpha)^*\in\bigcup_{m<n}\mathrm{Acc}_m$.
\end{lemma}

\begin{proof}
We focus on $\alpha$ with either $\{0,-1\}\subseteq\mathbb{FC}(\alpha)$ or $\mathbb{FC}(\alpha)=\{0\}$. (If $0\not\in\mathbb{FC}(\alpha)$, we are covered by the previous lemma. If $-1\not\in\mathbb{FC}(\alpha)$ but $\overline{\mathbb{FC}}(\vartheta\alpha)>-\infty$ then there is an $\alpha'$ with $\{0,-1\}\subseteq\mathbb{FC}(\alpha')$ and $(\vartheta\alpha')^*=(\vartheta\alpha)^*$.)

We proceed by induction on $\alpha\in\mathrm{Acc}_n$ with this property. Since $\mathbb{G}(\vartheta\alpha)<\mathbb{G}(\alpha)$, we certainly have $(\vartheta\alpha)^*\in\bigcup_{m<n}M_m$. Let $m=\mathbb{G}((\vartheta\alpha)^*)$.
  
  It suffices to show that for all $\gamma<(\vartheta\alpha)^*$ in $M_m$, $\gamma\in\mathrm{Acc}_m$. We show this by a side structural induction on $\gamma$.

  If $\gamma$ is $\Omega^{(0)}$, this is immediate since $\Omega^{(0)}$ is the smallest element of formal cardinality $0$. If $\gamma=\#\{\gamma_i\}$ then each $\gamma_i$ is either in $\mathrm{Acc}_k$ for some $k<m$, or $\gamma_i\in M_m$, $\gamma_i<(\vartheta\alpha)^*$, and therefore $\gamma_i\in\mathrm{Acc}_m$ by the side inductive hypothesis. Then, by closure under addition, $\gamma\in\mathrm{Acc}_m$.

  Similarly, if $\gamma=\omega^{\gamma'}$ then $\gamma'<(\vartheta\alpha)^*$, so $\gamma'\in\mathrm{Acc}_m$, and therefore $\omega^{\gamma'}\in\mathrm{Acc}_m$.

  For the main case, suppose $\gamma=\vartheta\gamma'$. If $\overline{\mathbb{FC}}(\gamma')<0$ then, by the previous lemma, $\gamma\in\mathrm{Acc}_m$. Otherwise, since $\vartheta\gamma'\in M_m$, we must have either $\overline{\mathbb{FC}}(\vartheta\gamma')=-\infty$ or $-1\in \mathbb{FC}(\gamma')$. If $\delta\in K^{< 0}\gamma'$ then either $\delta\in K^{< 0}\gamma$ (so $\delta^*\in\mathrm{Acc}_{\mathbb{G}(\delta^*)}$ since $\gamma\in M_m$) or $\delta\in K^{< 0}\gamma'\setminus K^{< 0}\gamma$, and therefore $\overline{\mathbb{FC}}(\delta)=-1$ and then we have $\delta<\vartheta\alpha$ and therefore, by the side inductive hypothesis, $\delta\in\mathrm{Acc}_m$. It follows that $\gamma'\in M_n$, and therefore $\vartheta\gamma'\in\mathrm{Acc}_m$ by the main inductive hypothesis.%
  
  Otherwise, there is some $\delta\in K^{< 0}\alpha$ with $\gamma\leq\delta$, and since $\delta\in\mathrm{Acc}_m$, also $\gamma\in\mathrm{Acc}_m$.
\end{proof}

\end{document}
